% GENERATED FILE. Edit canonical lessons, then run npm run generate:books.
% canonical-source-hash: fnv1a64:9e6df6f098354800
% canonical-lessons: RU-C271-vnezapno, RU-C271-mgnovenno, RU-C271-totchas, RU-C271-nemedlenno, RU-C271-mgnovenie

\chapter{Suddenly and at once}
\label{ch:ru-suddenly-and-at-once}
\hlchaptermodality{271}

\begin{chapteropening}
\textbf{By the end of this chapter:} \emph{I can say that something happened suddenly, in a flash, at once, without delay and in an instant.}

Complete the last lesson of chapter 271: I can say that something happened suddenly, in a flash, at once, without delay and in an instant.
\end{chapteropening}

\section[\texorpdfstring{\ru{внезапно} (vnezápno)}{внезапно (vnezápno)}]{\ru{внезапно} (vnezápno) — suddenly, unexpectedly}

\label{lesson:RU-C271-vnezapno}



\begin{quote}
Before the new one: say the Russian for a task, then the Russian for experience.
\end{quote}

\subsection*{You'll want to know: \ru{внезапно}}
\textbf{\ru{внезапно}} (\emph{vnezápno}) — \textquotedblleft{}suddenly, unexpectedly\textquotedblright{}.

\begin{quote}
\ru{Внезапно} \ru{в} \ru{комнате} \ru{стало} \ru{темно}.
\end{quote}

\begin{quote}
\emph{Vnezápno v kómnate stálo temnó.}
\end{quote}

\begin{quote}
Suddenly it went dark in the room.
\end{quote}

The adverb of \emph{\ru{внезапный}} (\emph{vnezápnyy}) \textquotedblleft{}sudden\textquotedblright{}. A little more bookish than \emph{\ru{вдруг}}, and common in stories.

\begin{cousinweb}
\emph{\ru{внезапно}} is the adverb of \emph{\ru{внезапный}} \textquotedblleft{}sudden\textquotedblright{}, built with the old prefix \emph{\ru{вне}-}. The origin of the stem \emph{-\ru{запн}-} is not settled, and EtymDB lists no parent.
\end{cousinweb}

\begin{grammarlens}[title={What you've built}]
One more word for telling, explaining and describing at B1.
\end{grammarlens}

\subsection*{Your turn}
\begin{itemize}
\raggedright
  \item \emph{Say it:} \emph{vnezápno}
  \item \emph{Say it:} \emph{vnezápno}, once more
  \item \emph{Say it:} a sentence of your own with \emph{vnezápno}
  \item \emph{Recall:} say the Russian for a task, then the Russian for experience, then say \emph{vnezápno} again
\end{itemize}

\begin{tcolorbox}[breakable,colback=teal!4,colframe=teal!35!black,title={Before you move on}]
What does \ru{внезапно} mean? (\textquotedblleft{}Suddenly, unexpectedly\textquotedblright{}.) Say it once more.
\end{tcolorbox}
\section[\texorpdfstring{\ru{мгновенно} (mgnovénno)}{мгновенно (mgnovénno)}]{\ru{мгновенно} (mgnovénno) — instantly, in a flash}

\label{lesson:RU-C271-mgnovenno}



\begin{quote}
Before the new one: say the Russian for experience, then the Russian for suddenly.
\end{quote}

\subsection*{You'll want to know: \ru{мгновенно}}
\textbf{\ru{мгновенно}} (\emph{mgnovénno}) — \textquotedblleft{}instantly, in a flash\textquotedblright{}.

\begin{quote}
\ru{Он} \ru{мгновенно} \ru{понял}, \ru{что} \ru{случилось}.
\end{quote}

\begin{quote}
\emph{On mgnovénno pónyal, shto sluchílos.}
\end{quote}

\begin{quote}
He instantly understood what had happened.
\end{quote}

The adverb of \emph{\ru{мгновенный}} (\emph{mgnovénnyy}) \textquotedblleft{}instant\textquotedblright{}. \emph{\ru{Мгновенно}} is faster still than \emph{\ru{быстро}}.

\begin{cousinweb}
\emph{\ru{мгновенно}} is built on \emph{\ru{мгновение}} \textquotedblleft{}an instant\textquotedblright{}, which is inherited through Old East Slavic \emph{mĭgnovenije} from Proto-Slavic \emph{mĭgnovenĭje}, built on \emph{mĭgnǫti} \textquotedblleft{}to blink\textquotedblright{} (EtymDB): as quick as a blink.
\end{cousinweb}

\begin{grammarlens}[title={What you've built}]
One more word for telling, explaining and describing at B1.
\end{grammarlens}

\subsection*{Your turn}
\begin{itemize}
\raggedright
  \item \emph{Say it:} \emph{mgnovénno}
  \item \emph{Say it:} \emph{mgnovénno}, once more
  \item \emph{Say it:} a sentence of your own with \emph{mgnovénno}
  \item \emph{Recall:} say the Russian for experience, then the Russian for suddenly, then say \emph{mgnovénno} again
\end{itemize}

\begin{tcolorbox}[breakable,colback=teal!4,colframe=teal!35!black,title={Before you move on}]
What does \ru{мгновенно} mean? (\textquotedblleft{}Instantly, in a flash\textquotedblright{}.) Say it once more.
\end{tcolorbox}
\section[\texorpdfstring{\ru{тотчас} (totchás)}{тотчас (totchás)}]{\ru{тотчас} (totchás) — at once, immediately}

\label{lesson:RU-C271-totchas}



\begin{quote}
Before the new one: say the Russian for suddenly, then the Russian for instantly.
\end{quote}

\subsection*{You'll want to know: \ru{тотчас}}
\textbf{\ru{тотчас}} (\emph{totchás}) — \textquotedblleft{}at once, immediately\textquotedblright{}.

\begin{quote}
\ru{Врач} \ru{тотчас} \ru{пришёл}.
\end{quote}

\begin{quote}
\emph{Vrach totchás prishyól.}
\end{quote}

\begin{quote}
The doctor came at once.
\end{quote}

Stressed on the end, \emph{totchás} (\emph{tótchas} is also heard). Bookish and a little old-fashioned; in speech Russians say \emph{\ru{сразу}}. Often strengthened: \emph{\ru{тотчас} \ru{же}}.

\begin{cousinweb}
\emph{\ru{тотчас}} joins \emph{\ru{тот}} \textquotedblleft{}that\textquotedblright{} and \emph{\ru{час}} \textquotedblleft{}hour\textquotedblright{}, which is inherited from Proto-Slavic \emph{časŭ} \textquotedblleft{}time\textquotedblright{} (EtymDB): that very hour.
\end{cousinweb}

\begin{grammarlens}[title={What you've built}]
One more word for telling, explaining and describing at B1.
\end{grammarlens}

\subsection*{Your turn}
\begin{itemize}
\raggedright
  \item \emph{Say it:} \emph{totchás}
  \item \emph{Say it:} \emph{totchás}, once more
  \item \emph{Say it:} a sentence of your own with \emph{totchás}
  \item \emph{Recall:} say the Russian for suddenly, then the Russian for instantly, then say \emph{totchás} again
\end{itemize}

\begin{tcolorbox}[breakable,colback=teal!4,colframe=teal!35!black,title={Before you move on}]
What does \ru{тотчас} mean? (\textquotedblleft{}At once, immediately\textquotedblright{}.) Say it once more.
\end{tcolorbox}
\section[\texorpdfstring{\ru{немедленно} (nemédlenno)}{немедленно (nemédlenno)}]{\ru{немедленно} (nemédlenno) — immediately, without delay}

\label{lesson:RU-C271-nemedlenno}



\begin{quote}
Before the new one: say the Russian for instantly, then the Russian for at once.
\end{quote}

\subsection*{You'll want to know: \ru{немедленно}}
\textbf{\ru{немедленно}} (\emph{nemédlenno}) — \textquotedblleft{}immediately, without delay\textquotedblright{}.

\begin{quote}
\ru{Он} \ru{немедленно} \ru{позвонил} \ru{домой}.
\end{quote}

\begin{quote}
\emph{On nemédlenno pozvoníl domóy.}
\end{quote}

\begin{quote}
He phoned home immediately.
\end{quote}

Literally \textquotedblleft{}not slowly\textquotedblright{}. Firmer than \emph{\ru{сразу}}; common in orders and official notices.

\begin{cousinweb}
\emph{\ru{немедленно}} is built on \emph{\ru{не}} \textquotedblleft{}not\textquotedblright{} and \emph{\ru{медленно}} \textquotedblleft{}slowly\textquotedblright{} (EtymDB lists \emph{\ru{медленно}} as its source), the adverb of \emph{\ru{медленный}} \textquotedblleft{}slow\textquotedblright{}: without lingering.
\end{cousinweb}

\begin{grammarlens}[title={What you've built}]
One more word for telling, explaining and describing at B1.
\end{grammarlens}

\subsection*{Your turn}
\begin{itemize}
\raggedright
  \item \emph{Say it:} \emph{nemédlenno}
  \item \emph{Say it:} \emph{nemédlenno}, once more
  \item \emph{Say it:} a sentence of your own with \emph{nemédlenno}
  \item \emph{Recall:} say the Russian for instantly, then the Russian for at once, then say \emph{nemédlenno} again
\end{itemize}

\begin{tcolorbox}[breakable,colback=teal!4,colframe=teal!35!black,title={Before you move on}]
What does \ru{немедленно} mean? (\textquotedblleft{}Immediately, without delay\textquotedblright{}.) Say it once more.
\end{tcolorbox}
\section[\texorpdfstring{\ru{мгновение} (mgnovéniye)}{мгновение (mgnovéniye)}]{\ru{мгновение} (mgnovéniye) — an instant, a moment}

\label{lesson:RU-C271-mgnovenie}



\begin{quote}
Before the new one: say the Russian for at once, then the Russian for immediately.
\end{quote}

\subsection*{You'll want to know: \ru{мгновение}}
\textbf{\ru{мгновение}} (\emph{mgnovéniye}) — \textquotedblleft{}an instant, a moment\textquotedblright{}.

\begin{quote}
\ru{На} \ru{одно} \ru{мгновение} \ru{все} \ru{замолчали}.
\end{quote}

\begin{quote}
\emph{Na odnó mgnovéniye vse zamolcháli.}
\end{quote}

\begin{quote}
For an instant everyone fell silent.
\end{quote}

A neuter noun. \emph{\ru{В} \ru{одно} \ru{мгновение}} — in an instant; \emph{\ru{в} \ru{последнее} \ru{мгновение}} — at the last moment.

\begin{cousinweb}
\emph{\ru{мгновение}} is inherited through Old East Slavic \emph{mĭgnovenije} from Proto-Slavic \emph{mĭgnovenĭje}, built on \emph{mĭgnǫti} \textquotedblleft{}to blink\textquotedblright{}, a cousin of \emph{\ru{мигать}} \textquotedblleft{}to blink\textquotedblright{} (EtymDB).
\end{cousinweb}

\begin{grammarlens}[title={What you've built}]
One more word for telling, explaining and describing at B1.
\end{grammarlens}

\subsection*{Your turn}
\begin{itemize}
\raggedright
  \item \emph{Say it:} \emph{mgnovéniye}
  \item \emph{Say it:} \emph{mgnovéniye}, once more
  \item \emph{Say it:} a sentence of your own with \emph{mgnovéniye}
  \item \emph{Recall:} say the Russian for at once, then the Russian for immediately, then say \emph{mgnovéniye} again
\end{itemize}

\begin{tcolorbox}[breakable,colback=teal!4,colframe=teal!35!black,title={Before you move on}]
What does \ru{мгновение} mean? (\textquotedblleft{}An instant, a moment\textquotedblright{}.) Say it once more.
\end{tcolorbox}
